\begin{proposition}[\texorpdfstring{$\delta_q$}{deltaq} 的黄金比例包络]\label{prop:pom-deltaq-golden-envelope}
\leanverified{paper\_pom\_deltaq\_golden\_envelope}
沿用注 \ref{rem:pom-collision-kernel-rh-phase} 的记号：对每个整数 \(q\ge 2\)，令 \(r_q\) 为碰撞核 \(A_q\) 的 Perron 根，并令
\(
\Lambda_q:=\max_{\mu\in\sigma(A_q)\setminus\{r_q\}}|\mu|
\)
与
\(
\delta_q:=\log(\Lambda_q^2/r_q)
\)。
则对每个 \(q\ge 2\) 恒有严格上界 \(\delta_q<\log r_q\)。
并且在 Fold\(_m\) 的碰撞矩 \(S_q(m)=\sum_x d_m(x)^q\) 框架下，对每个 \(q\ge 1\) 有
\[
\log r_q\le \Bigl(\frac{q}{2}+1\Bigr)\log\varphi,
\]
从而 \(\delta_q\le (\frac{q}{2}+1)\log\varphi\)。
\end{proposition}
\begin{proof}
由于 \(\Lambda_q<r_q\)，有
\(
\delta_q=\log(\Lambda_q^2/r_q)<\log(r_q^2/r_q)=\log r_q
\)。
\par
另一方面，对任意 \(m\) 由平凡界 \(S_q(m)\le |X_m|\cdot D_m^q\)（其中 \(D_m:=\max_{x\in X_m}d_m(x)\)）得
\[
S_q(m)^{1/m}\le |X_m|^{1/m}\,D_m^{q/m}.
\]
令 \(m\to\infty\)。由命题 \ref{prop:pom-projection-entropy} 知 \(|X_m|=F_{m+2}=\Theta(\varphi^m)\)，且由定理 \ref{thm:pom-max-fiber} 知 \(D_m=\Theta(\varphi^{m/2})\)，故
\(
\lim_{m\to\infty}|X_m|^{1/m}=\varphi
\)
与
\(
\lim_{m\to\infty}D_m^{1/m}=\varphi^{1/2}
\)。
于是得到
\(
r_q=\lim_{m\to\infty}S_q(m)^{1/m}\le \varphi\cdot(\varphi^{1/2})^{q}=\varphi^{q/2+1}
\)，即所示上界。证毕。
\end{proof}

\begin{corollary}[近线性拟合外推的显式破缺阈值]\label{cor:pom-deltaq-linear-extrapolation-break}
\leanverified{paper\_pom\_deltaq\_linear\_extrapolation\_break}
若将式 \eqref{eq:fold-collision-delta-linear-fit-q18-23} 的近线性律外推为对一切 \(q\ge 18\) 成立的全局下界
\[
\delta_q\ge 0.263694\,q-0.842355\qquad(\forall q\ge 18),
\]
则该下界与命题 \ref{prop:pom-deltaq-golden-envelope} 的 \(\delta_q\le (\frac{q}{2}+1)\log\varphi\) 在 \(q\ge 58\) 处矛盾。
因此该线性律不可能在 \(q\) 方向无限延拓，且必存在
\[
q_0\in\{18,19,\dots,57\}
\]
使得 \(\delta_{q_0}<0.263694\,q_0-0.842355\)。
\end{corollary}
\begin{proof}
由命题 \ref{prop:pom-deltaq-golden-envelope} 得对所有 \(q\ge 18\) 有
\[
0.263694\,q-0.842355\le \delta_q\le \Bigl(\frac{q}{2}+1\Bigr)\log\varphi.
\]
解不等式
\(
0.263694\,q-0.842355\le \frac12(\log\varphi)\,q+\log\varphi
\)
得 \(q\le 57.33\ldots\)，故当 \(q\ge 58\) 必矛盾，从而在集合 \(\{18,\dots,57\}\) 内至少有一处破缺。证毕。
\end{proof}
